\section{Method}
\label{sec:method}

Our pipeline converts thermal-infrared video into event streams through four
stages (Fig.~\ref{fig:pipeline}): (i) radiometric-to-display audit and AGC
de-flickering, (ii) temporal interpolation, (iii) threshold-crossing event
generation with hardware-specific noise models, and (iv) packaging with
label alignment. Each stage exposes thermal-specific parameters that we
study systematically in Sec.~\ref{sec:fidelity}.

\subsection{AGC de-flickering}
\label{sec:agc}
Uncooled thermal cameras apply automatic gain control (AGC) that continuously
re-maps the radiometric signal to the display range. Frame-to-frame AGC drift
and abrupt re-ranging events create \emph{global} photometric flicker that, fed
into a threshold-crossing event simulator, produces scene-wide spurious events
that drown genuine motion signal. This effect is documented in real thermal
video by ``Thermal is Always Wild''~\cite{thermalwild2026} and makes AGC
handling a contribution rather than an engineering detail.

We model per-frame AGC as a scalar gain about a photometric pivot,
$\mathrm{disp} = g\,(q - \mathrm{piv}) + \mathrm{piv}$, where the pivot
(histogram middle for typical thermal AGC) is pinned. We estimate
$\mathrm{piv}$ as the temporal median of per-frame medians, fit $g$ per frame
by least squares on pivot-relative deviations of a dense percentile profile
(21 levels), and invert. Single-parameter pivot-centred correction keeps
bulk near-pivot pixels---which dominate real scenes---pinned even under
gain-estimation noise, avoiding correction-induced event floods that naive
two-parameter affine correction produces (Sec.~\ref{sec:fidelity}). The
correction is gated off when no significant gain variation is detected. On
synthetic scenes with known ground truth, our de-flicker recovers the
radiometric event-count floor \emph{exactly}.

\subsection{Temporal interpolation}
\label{sec:interp}
Threshold-crossing simulators require high temporal sampling; thermal sources
range from \SI{7.5}{\hertz} (HIT-UAV) to \SI{30}{\hertz} (FLIR). We compare
linear blending, cubic (Catmull--Rom), and optical-flow-guided bidirectional
warping (Farneb\"ack~\cite{farneback2003}) under held-out reconstruction.
Contrary to visible-light practice, learned or flow-guided interpolation does
\emph{not} help on thermal video at these rates (Sec.~\ref{sec:e1}); we
therefore default to linear interpolation, an E1 finding of independent
interest.

\subsection{Event generation with thermal physics}
\label{sec:sim}
Our simulator implements three modes on a shared backend:
(i) \textbf{v2e}---classic threshold crossing with per-pixel threshold
mismatch, log-domain shot noise, refractory period and leak
events~\cite{hu2021v2e}; (ii) \textbf{volt}---DVS-Voltmeter-style Brownian
threshold random walk~\cite{lin2022voltmeter}; and (iii)
\textbf{poisson}---an uncooled Poisson-bolometer model in which event
emission is an inhomogeneous Poisson process whose rate grows exponentially
with log-intensity overdrive, emulating spintronic stochastic-switching LWIR
pixels~\cite{mousa2026poisson}. Thermal physics enters through (a) an
optional first-order low-pass on log intensity emulating the microbolometer
time constant (ADV2E-style~\cite{jiang2024adv2e}), and (b) a
threshold$\leftrightarrow$NETD mapping obtained by calibrating event yield
inside known object boxes (Sec.~\ref{sec:calibration}).

\subsection{Output packaging}
Events are stored as compressed $(t,x,y,p)$ streams (float32 timestamps),
with labels aligned to source frame times by nearest-frame association
following the Prophesee protocol~\cite{propheseetoolbox}.
